\documentclass[tikz,border=4pt]{standalone}
\usepackage{ctex}
\usepackage{amsmath,amssymb}
\usepackage{pgfplots}
\pgfplotsset{compat=1.18}
\usepgfplotslibrary{groupplots,fillbetween}
\begin{document}
\begin{tikzpicture}
\begin{groupplot}[
  group style={
    group size=3 by 1,
    horizontal sep=1.5cm,
  },
  width=5.5cm, height=5.0cm,
  xlabel style={font=\small},
  ylabel style={font=\small},
  tick label style={font=\footnotesize},
  grid=both, grid style={gray!20},
  axis lines=left,
  samples=200,
]

% Plot 1: PDF and its area property
\nextgroupplot[
  title={PDF：$f(x)\geq 0$，$\int f=1$},
  title style={font=\small},
  xmin=-4, xmax=4,
  ymin=-0.05, ymax=0.55,
  xlabel={$x$}, ylabel={$f(x)$},
]
% Gaussian PDF
\addplot[name path=gauss, blue, very thick, domain=-4:4]
  {exp(-x^2/2)/sqrt(2*3.14159)};
\addlegendentry{$\mathcal{N}(0,1)$}
% Fill under curve
\addplot[name path=zero, draw=none, domain=-4:4] {0};
\addplot[blue!20, opacity=0.5] fill between [of=gauss and zero];
% Area annotation
\node[font=\footnotesize, fill=white] at (axis cs:0, 0.25) {$\int_{-\infty}^{\infty}f\,dx=1$};
% Height annotation
\draw[->, gray] (axis cs:0.1,0.4) -- (axis cs:0.1, 0.398) node[above, font=\footnotesize]{最高点};
% f(x) can exceed 1 note
\node[font=\footnotesize, fill=yellow!20, draw=orange, rounded corners,
      align=center] at (axis cs:2.5, 0.4) {$f(x)$ 可 $>1$\\不是概率！};

% Plot 2: CDF as integral of PDF
\nextgroupplot[
  title={CDF：$F(x)=\int_{-\infty}^{x}f\,dt$},
  title style={font=\small},
  xmin=-4, xmax=4,
  ymin=-0.05, ymax=1.15,
  xlabel={$x$}, ylabel={$F(x)$},
]
\addplot[red, very thick, domain=-4:4]
  {1/(1+exp(-1.7*x))};
\addlegendentry{$F(x)$（逻辑近似）}
% annotation at 0
\addplot[mark=*, mark size=2pt, red] coordinates {(0, 0.5)};
\node[font=\footnotesize, right] at (axis cs:0.1, 0.5) {$F(0)=0.5$};
% asymptotes
\addplot[gray, dashed, domain=-4:4] {1};
\addplot[gray, dashed, domain=-4:4] {0};
\node[font=\footnotesize, right] at (axis cs:2.2, 1.05) {$\lim_{x\to\infty}F=1$};
\node[font=\footnotesize, right] at (axis cs:2.2, 0.07) {$\lim_{x\to-\infty}F=0$};
% derivative annotation
\node[font=\footnotesize, fill=white] at (axis cs:-2, 0.7) {$F'(x)=f(x)$};

% Plot 3: Expected value geometric
\nextgroupplot[
  title={期望：$E[X]=\int xf(x)\,dx$},
  title style={font=\small},
  xmin=-4, xmax=6,
  ymin=-0.05, ymax=0.55,
  xlabel={$x$}, ylabel={$f(x)$},
]
% Skewed distribution: mixture or shifted Gaussian
\addplot[name path=pdf2, blue, very thick, domain=-4:6]
  {0.6*exp(-(x-1)^2/0.8)/sqrt(0.8*3.14159) + 0.4*exp(-(x-3)^2/1.2)/sqrt(1.2*3.14159)};
\addlegendentry{双峰混合分布}
\addplot[name path=zero2, draw=none, domain=-4:6] {0};
\addplot[blue!15] fill between [of=pdf2 and zero2];
% E[X] vertical line ~ 1.8 (weighted average of 1 and 3)
\addplot[red, very thick, dashed] coordinates {(1.8, 0)(1.8, 0.4)};
\node[font=\footnotesize, red, right] at (axis cs:1.85, 0.3) {$E[X]\approx 1.8$};
\node[font=\footnotesize, fill=white] at (axis cs:-1.5, 0.4) {"重心"位置};
% modes
\addplot[mark=triangle*, mark size=3pt, orange] coordinates {(1, 0.38)};
\node[font=\footnotesize, above] at (axis cs:1, 0.38) {众数$_1$};
\addplot[mark=triangle*, mark size=3pt, orange] coordinates {(3, 0.23)};
\node[font=\footnotesize, above] at (axis cs:3, 0.23) {众数$_2$};

\end{groupplot}

\node[font=\normalsize\bfseries] at (9.5, 5.7) {PDF / CDF / 期望的几何意义};
\end{tikzpicture}
\end{document}
